\documentclass[11pt]{article}
\usepackage[margin=1in]{geometry}
\usepackage{amsmath}

\title{Paper 1 (SAR-CRP v2 Core) --- MVP Results Note}
\author{}
\date{\today}

\begin{document}
\maketitle

\section{Scope}

This note reports the Minimum Viable Product (MVP) experiment defined in
\emph{SAR-CRP v2 FINAL READY Proposal}, \S33: three methods (Static,
Full Reoptimization, SAR-CRP Core) on one small yard layout
(\texttt{mvp\_small\_01}, 4 stacks, 10 containers), under a synthetic
dynamic event stream (order swap, urgent insertion, ETA shift,
probability update, stale information), 10 random seeds per method per
uncertainty level, swept over \texttt{low}/\texttt{medium}/\texttt{high}
uncertainty (spec \S23 Experiment 1 factor). All computation ran on the
remote CPU server (\texttt{a100-B}, \texttt{story\_research} conda env);
the laptop was used only to edit code.

\textbf{Instance note.} The initial retrieval order was first set to match
each stack's natural top-to-bottom ejection order, which made the base plan
require zero relocations --- exactly the ``benchmark too easy'' failure
mode spec \S20 SC1 warns against. The order was corrected to force real
blocking (7 relocations in the base plan) before the results below were
produced.

\textbf{Correctness fix.} An earlier version of this experiment showed
SAR-CRP with unexplained cost spikes on episodes where it never actually
replanned (\texttt{fallback\_rate}=1.0). Root cause: \texttt{minimal\_}
\texttt{feasibility\_repair} and three of the local-search neighborhood
operators (N3/N4/N5) reused \texttt{plan\_old}'s own \texttt{Action}
objects while building a repair candidate, then mutated their
\texttt{step\_index} in place. Since \texttt{sarcrp\_core.replan} returns
\texttt{plan\_old} \emph{by reference} on every KEEP decision, and the
simulator keeps reusing that same object for the rest of the episode,
this silently corrupted the indices that \texttt{stability\_cost} and
\texttt{data\_confidence\_cost} key off of -- even when the plan was never
actually updated. All four sites now deep-copy kept actions before
renumbering. The numbers below are from the corrected code.

\section{Results}

\begin{table}[h]
\centering
\small
\begin{tabular}{lrrrrr}
\hline
Uncertainty & Method & Total cost & Op.\ cost & Changed actions & Fallback rate \\
            &        & (mean)     & (mean)    & (mean)           &               \\
\hline
low    & Static              & 7.034  & 7.034 & 0.0  & 1.000 \\
low    & Full Reoptimization & 19.464 & 6.803 & 55.6 & 0.000 \\
low    & SAR-CRP Core        & 7.034  & 7.034 & 0.0  & 1.000 \\
\hline
medium & Static              & 7.032  & 7.032 & 0.0  & 1.000 \\
medium & Full Reoptimization & 20.210 & 6.510 & 66.0 & 0.000 \\
medium & SAR-CRP Core        & 7.032  & 7.032 & 0.0  & 1.000 \\
\hline
high   & Static              & 7.040  & 7.040 & 0.0  & 1.000 \\
high   & Full Reoptimization & 21.433 & 6.334 & 68.8 & 0.000 \\
high   & SAR-CRP Core        & 7.040  & 7.040 & 0.0  & 1.000 \\
\hline
\end{tabular}
\caption{MVP results averaged over 10 seeds per row. Source: \texttt{paper1\_sarcrp/experiments/results/mvp\_results.csv}.}
\end{table}

\textbf{SAR-CRP ties Static exactly, at every uncertainty level, across
all 30 seeds.} This is a real and explainable result, not a bug: the
trigger did fire at least once in several episodes (e.g.\ seed 7,
medium, event 5, an \texttt{ORDER\_SWAP} with impact 0.361 $\ge
\theta_{impact}=0.30$) but the outer decision still resolved to KEEP,
because $J(P_{old})=J(P_{best})$ exactly. The reason is structural, not a
tuning artifact: spec \S11.2's $\mathit{RetrievalDelay}(P)$ term only
sums position over \emph{urgent} containers, and only
\texttt{URGENT\_INSERTION} events populate the urgent set. A pure
\texttt{ORDER\_SWAP}/\texttt{ETA\_shift}/\texttt{PROBABILITY\_UPDATE} event
changes the queue but changes $C_{op}(P)$ for \emph{no} plan whatsoever
(reordering has zero effect on relocation count or on a delay term that
only counts non-urgent... i.e.\ empty containers), so there is nothing for
any repair candidate to improve on. In this seed set, the
\texttt{URGENT\_INSERTION} events that did occur never happened to also
cross $\theta_{impact}$.

Full Reoptimization tells the complementary half of the story: it
achieves a genuinely \emph{lower} operational cost than Static/SAR-CRP
(6.3--6.8 vs.\ 7.03) by re-solving on every single event, but pays for
that with 56--69 changed actions per episode, which under $\lambda=1.0$
dominates its total cost (19--21 vs.\ 7.03). SAR-CRP correctly recognizes
that this trade is not worth it and defaults to Static's behavior instead
of Full Reoptimization's churn -- which is the intended thesis of Paper 1
(stability-aware trade-off), just not yet demonstrated by a scenario where
repairing is unambiguously worthwhile.

\section{Decision gate (spec \S33)}

\begin{itemize}
  \item SAR-CRP total cost $<$ Static total cost: \textbf{TIE, not a strict PASS} at all three levels (equal, not lower)
  \item SAR-CRP changed-actions (stability proxy) $<$ Full Reoptimization: \textbf{PASS} at all three levels (0.0 $\ll$ 55.6--68.8)
  \item SAR-CRP operational cost $\le$ 1.20 $\times$ Full Reoptimization's operational cost: \textbf{PASS} at all three levels (7.03--7.04 $\le$ 1.2 $\times$ 6.3--6.8)
\end{itemize}

The operational-cost condition is now computed from $C_{op}(P)$ alone
(spec \S11) rather than blended total cost, and is upper-bound-only:
SAR-CRP sitting \emph{below} Full Reoptimization's operational cost is
never treated as a failure, only sitting meaningfully above it is. This
replaces the earlier $\pm20\%$ band on total cost, which flagged SAR-CRP's
better-than-Full-Reopt result as a failure purely because the proxy
conflated operational and stability cost.

\section{Next step}

Two of three conditions hold cleanly; the first is a tie rather than a
genuine win, and the diagnosis above shows why: this seed set never
produced an urgent-container scenario with high enough impact to make
repairing pay off under the current cost proxy. Before proceeding to the
full Implementation Roadmap Phase 6 (\S27) --- the 6-baseline,
5-experiment Q1 suite with the \S23.6 statistical protocol (20+ seeds,
Wilcoxon signed-rank, Holm--Bonferroni) --- the following should be
addressed:
\begin{itemize}
  \item Construct at least one scenario with a clear, high-impact urgent
    insertion (e.g.\ force one early in the event stream) to verify
    SAR-CRP can demonstrate a real win, not just parity with Static, before
    trusting the full sweep.
  \item Re-examine whether spec \S11.2's urgent-only $\mathit{RetrievalDelay}$
    term is too narrow for the MVP's synthetic benchmark, or whether this
    is simply what a deliberately simple proxy should look like --- this is
    a modeling decision, not a code defect, and belongs in the paper's
    Limitations discussion either way.
  \item N3/N5 local-search neighborhoods are now implemented (spec \S15.1,
    \S46.2); the CRP\_RL-tail candidate's \texttt{forbidden\_moves}/
    \texttt{max\_changed\_actions} constraint forwarding remains deferred.
  \item Real code and benchmark instances for two of the three verified
    base papers were located during this iteration (Shin et al.\ 2026's
    CRP\_RL: \url{https://github.com/operagang/CRP_RL}, MIT-licensed code +
    CC-BY-4.0 data, including pretrained checkpoints; Zhou \& Zhang 2024:
    \url{https://github.com/zhou-sf/rt-scrp}, MATLAB). Swapping the greedy
    surrogate solver for the real trained CRP\_RL model, and/or adopting
    the Lee/Shin benchmark instances in place of the hand-built MVP
    instance, is a separate scoped follow-up (schema adapter required --
    their environment uses a bay/row/tier tensor representation and a
    real travel-time cost model, not this MVP's flat-stack schema).
\end{itemize}

\end{document}
